\phantomsection\label{app:probing-proof}
\begin{proof}
Fix a summand. Write $wz_i$ for its outstanding bits on commodity $i$.
Record the bit positions in order at each endpoint, so any subsequent
identity on these positions is valid under the surrounding contractions.
Let $S$ count the normalized capacity fees and linear costs already
incurred, less a credit of $2M/(nw)$ for each external catalyst bit
used in a replacement. An unused catalyst bit can be routed through
its column center for exactly its credit. Thus, once all commodity
demands are met, the total cost is at most $S+2MN_c/(nw)$. We control
the cost of unfinished demands by the potential $F_s(z)$.

\paragraph{Marginal cost of a demand.}
Let $\mu$ be the strong convexity constant of $f$. At the maximizer
$x=x_s(z)$ the box optimality condition gives
\[
                        z_i\ge s\nabla_i f(x)\ge s\mu x_i .
\]
For the second inequality, vary only coordinate $i$ and use the
vanishing of $\nabla_i f$ on its lower face. On that face the first
inequality follows from $z_i\ge0$. In particular, a price bounded away
from zero guarantees stock when $s$ is bounded away from zero. Choose
an integer $B$ greater than every gradient coordinate bound. For
$z_i\ge B$ we have $x_i=a_i$. We can therefore fulfill all stock above
$B$ along the prescribed route at cost at most $a_i/n$ per unit while
decreasing $F_s$ by precisely that amount per unit. We apply this cap
after every stock increase.

With the inventory norm induced by $\langle\cdot,\cdot\rangle$, the
two optimality inequalities and strong convexity give
\[
 \|x_s(z)-x_{s'}(z')\|
       \le \frac{\|z-z'\|+C|s-s'|}{\mu\min(s,s')}.
\]
Comparing the two maxima at their respective optimizers now gives
$\nabla_z F_s(z)=x_s(z)$ and $\partial_s F_s(z)=-f(x_s(z))$.
For bounded inventories and $s,s'$ in a compact subinterval of
$(0,1]$, the preceding estimate makes the first-order remainder
uniformly quadratic. We will use these derivatives both to price a
local replacement and to pay for new probes.

\paragraph{A failed test gives a cheaper exact identity.}
On a ball that fails the test, demand one bit of each positive-priced
commodity and of every catalyst pair. Let $L$ be the sum of their
assigned prices. The zero wire metric fails the test, so $L>0$.
Available replacements are commodities at prices $x_i$, catalyst pairs
at $2M$, and core links and currently tested probes at their lengths
$\lambda_e$. A replacement commodity enters the inventory. A
replacement catalyst pair can be installed immediately along its center
route.

Increase all replacement prices by the same sufficiently small amount
to obtain positive lengths $l_e$. On a nonsingleton ball, every
normalized spectral wire metric $d$ then satisfies
\[
             \sum_{e\ \mathrm{demanded}}d_e
               < L\max_{e\ \mathrm{available}}d_e/l_e .
\]
The connected core makes the maximum positive. If no increase works, take a
sequence of violations as the price increase tends to zero and rescale
each violating metric to make its maximum ratio one. The lengths of
a spanning tree bound the scales. In the limit every replacement
respects its original price bound, yet the sum of the demand distances
is at least $L$. Each demand must therefore attain its own bound.
This realizes the prices of all catalysts and positive-priced commodities,
and the zero-priced commodities have distance zero. The resulting metric
passes the test, a contradiction. On a singleton all demands can be
supplied locally.

By compactness the displayed inequality holds with some $L_-<L$ in
place of $L$. Take the product of width-$b_0$ demand wires. As supply,
take the direct sum over replacement options $e$ of a single wire of
width $\lfloor b_0 L_-/l_e\rfloor$ on $e$. For large fixed $b_0$ every
spectral point values the supply strictly above the demand, by the
uniform margin. Lemma~\ref{lem:spectral-comparison} supplies the $r$th
demand power from a direct sum of at most $C(b_0)$ copies of the $r$th
supply power. Expansion gives an ordinary sum over $F$ of restrictions
of replacement products. With $b=b_0r$ bits per demand and $m$
replacement options, each term costs at most $bL_-$, and there are at
most $C(b_0)m^r$ terms. Choosing first $b_0$ and then $r$ makes the
number of bits specifying a term arbitrarily small relative to $b$.

For $J=O(w)$ simultaneous copies, the outer selector records only the
histogram of the $J$ choices. It uses $O(\log(2w))$ bits because the
alphabet is fixed. Inside that summand, transmit the full choice word
on a core spanning tree, check its histogram, and sum over the words with
that histogram. The histogram fixes the installed widths, while each word
specifies offsets within them. In this way every original word
contributes its tensor with its original coefficient. The tree has
bounded total length, so the rate just obtained makes the cost of
transmitting the word less than half the saving $Jb(L-L_-)$. Fixed
tables, prefix sums and sorting of bit positions implement the local
operations with $O(w\,\mathrm{polylog}(h))$ trivalent gates since
$w\le h^A$. Each auxiliary deterministic check has a unique valid
history.

Suppose the mean discrepancy exceeds a fixed tolerance. Bounded overlap
and bounded discrepancies yield $\Omega(n)$ disjoint balls failing by
a fixed amount. In each bounded local format the vectors whose
discrepancy reaches this fixed lower bound form a compact set. At any one
vector the saving from its identity persists on a neighborhood where every selected
positive-priced commodity stays bounded away from zero. A finite
subcover gives a fixed menu with uniform margins. An $o(1)$
normalized-norm error in the grid approximation affects only $o(n)$
balls by more than a fixed coordinate tolerance.

Use $J=\lfloor\alpha w/b\rfloor$ copies on each selected ball, checking
the actual stock and skipping deficient balls. For fixed sufficiently
small $\alpha>0$ the stock inequality at the beginning of the proof
guarantees adequate inventory on every accurate ball while $s$ stays
bounded away from zero. The $o(n)$ inaccurate balls have bounded
cost per ball in units of $w$. In this simultaneous round, use the marginal prices $x_i$ for newly
created and fulfilled demands. Include the linear link costs and deduct
the $2M$ credit for each consumed external catalyst bit. After division by $nw$ these changes decrease the
first-order estimate for $S+F_s(z)$ by at least $c\alpha$, uniformly
for large $h$. The change in inventory has normalized norm $O(\alpha)$,
so its quadratic error is $O(\alpha^2)$. Fixing $\alpha$ sufficiently
small leaves a uniform positive decrease in every new summand. This
applies to both the core test and the test with one purchased activation.

\paragraph{Paying for probes when the core test passes.}
Fix rational $s$. Suppose the excess of $x=x_s(z)$ exceeds $1/4$ and
its core discrepancy is sufficiently small. Choose the activation
from the grid approximation to $x$. For some fixed $\zeta>0$ and
test tolerance, every $x'$ of at most this discrepancy with the
purchased probes satisfies, for all sufficiently large $h$,
\[
                      (s/2)\mathcal D_f(x',x)\ge P+\zeta .
\]
If such constants did not exist, a sequence with discrepancies tending
to zero and margins tending to zero or less would contradict
\eqref{eq:protocol-activation}. Bounded coordinates and gradients,
the Lipschitz gradient, and the normalized-norm approximation let us
use the actual maximizer in the inequality.

Buy the probes once. Hold half the inventory in reserve and start with
$z/2$ at time $s/2$. The potential is then $F_{s/2}(z/2)=F_s(z)/2$.
Feeding the reserve back at rate $z/s$ as time advances to $s$ would
keep $x$ fixed and restore the original potential in the absence of
replacements. Feed at this same rate on a fine fixed rational mesh.
Before each increment, apply rounds of simultaneous local replacements
until the current maximizer $x'$ has small discrepancy with the probes.
Along the feed the instantaneous increase in potential is at most
\[
 \langle z/s,x'\rangle-f(x')
 \le F_s(z)/s-\mathcal D_f(x',x)
 \le F_s(z)/s-2(P+\zeta)/s .
\]
The first inequality is the box optimality condition
$\langle z/s-\nabla f(x),x'-x\rangle\le0$ at the original maximizer.
Over the interval of length $s/2$, the feed therefore recovers at
most $F_s(z)/2-P-\zeta$. This pays the capacity fee $P$ and decreases
$S+F_s$ by $\zeta$. The sum of the quadratic mesh errors is proportional
to the mesh spacing. Rounding the split and the increments to whole bits,
returning the original reserve exactly, costs $o(1)$. A fine enough
mesh leaves a saving of at least $\zeta/2$.

The number of replacement rounds in a purchase can be bounded uniformly
even when $S$ is negative. The stock cap and bounded derivatives give
a bound for the total possible increase from the purchase, feed and
time changes without the test. Once the accumulated replacement savings
exceed that bound plus the required saving, complete the feed directly.

\paragraph{Finishing the payloads within budget.}
Start at $z=0$, $s=1$ and pay $P_{\mathrm{core}}\le1$. On a fine fixed
mesh, inject a total of exactly $w$ bits on each payload commodity.
Next decrease time to a fixed rational $s_0>0$ with $s_0f(a)<1$
uniformly. Before each increment, bring the excess down to at most
$1/2$. Apply a replacement round at large discrepancy, or buy probes
at small discrepancy when the excess exceeds $1/4$. The overlap of
these excess thresholds allows fixed-precision decisions. Purchased
links stay paid. To choose a new activation we test the core alone.
During the corresponding feed we test the core with its new probes.

At excess at most $1/2$, the mean payload price is at most $3/2$
during injection. During the time decrease, $f(x_s(z))\le j(s)+1/2$.
Since $\int_0^1 j(s)\,ds\le1$, each phase costs at most two on a fine
enough fixed mesh, including rounding. Finally route the remaining
requests along their prescribed core routes. Substituting $a$ in the
maximum defining $F_{s_0}(z)$ shows that this costs at most
$F_{s_0}(z)+s_0f(a)$, so the final $S$ is less than $10$.

As in the purchase step, the outer loop has a uniform stopping bound.
Even without the low-excess tests, the bounded derivatives control the
sum of all injection and time-change increases, and routing capped
stock directly has bounded cost. If the accumulated decreases exceed
this crude bound, complete the scheduled injections, time changes and
routing directly. Each trajectory thus performs uniformly bounded
numbers of reductions, purchases and replacement rounds.

Reserve $O(w)$ external bits per catalyst pair in advance and return
unused bits on their center routes. Aggregate repeated uses of each
available link into a single installed width with one capacity fee.
Bounded incidence and bounded routes give a bounded number of pieces
per site and step. Widths, boundaries and histograms require
$\mathrm{polylog}(h)$ bits per piece, with $O(h)$ more bits for site
identifiers. In the fixed local topology, pad each list to its maximum
length and sort position tags to split, concatenate or permute its
pieces. This takes $O(w\,\mathrm{polylog}(h))$ gates per site. Given
the layout, value and activation routines, the control handles the
fixed-precision menus, greedy disjoint-ball packings, stock checks,
histograms and link-width lists in $O(n\,\mathrm{polylog}(h))$ word
operations. Tolerance decisions have fixed slack, so their constants
may be hardwired. Mesh arithmetic uses rationals and rounded bit counts.

Use a unique encoding for each histogram and set unused selector fields
to zero. Each replacement is an exact identity under every contraction
on its ordered bit positions. Expanding the successive replacements
therefore gives exactly the payload matching times the external catalyst
identities with the prescribed multiplicities. The normalized cost is
at most $10+2MN_c/(nw)$ in every summand, as claimed.
\end{proof}
